\documentclass{beamer}
\usepackage{hyperref}
\usepackage[russian]{babel}
\usepackage{fontspec}

\usepackage{latexsym,amsmath,xcolor,multicol,booktabs,calligra}
\usepackage{graphicx,listings,stackengine}
\usepackage{array}

\setmainfont{CMU Serif}
\setsansfont{CMU Sans Serif}
\setmonofont{CMU Typewriter Text}

\author{Пучкин Станислав (СПбГУ)}
\title{Лямбды в C++}
\subtitle{Замыкания, захват, generic lambdas, type erasure}
\institute{C++ Семинар}
\date{\today}
\usepackage{CNU}

\AtBeginSection[]{}
\AtBeginSubsection[]{}

\definecolor{deepblue}{rgb}{0,0,0.5}
\definecolor{deepred}{rgb}{0.6,0,0}
\definecolor{deepgreen}{rgb}{0,0.5,0}
\definecolor{halfgray}{gray}{0.55}

\lstset{
    basicstyle=\ttfamily\scriptsize,
    keywordstyle=\bfseries\color{deepblue},
    emphstyle=\ttfamily\color{deepred},
    stringstyle=\color{deepgreen},
    commentstyle=\itshape\color{deepred},
    breaklines=true,
    columns=fixed,
    keepspaces=true,
    numbers=left,
    numberstyle=\tiny\color{halfgray},
    rulesepcolor=\color{red!20!green!20!blue!20},
    frame=shadowbox,
    language=C++
}

\begin{document}

\begin{frame}
    \titlepage
\end{frame}

% ----------------------------------------------------------------
\begin{frame}[fragile]{Лямбда — объект замыкания}
\begin{lstlisting}
// Без захвата: анонимная структура + приведение к func*
auto adder = [](int x, int y) { return x + y; };

// Компилятор генерирует примерно:
struct Closure {
    static int func(int x, int y) { return x + y; }
    operator decltype(&func)() const { return func; }
};
// Поэтому работает неявное приведение:
int (*pf)(int, int) = adder;      // ok

// Тип уникален — переприсвоить нельзя:
auto f = []{};
f = []{};  // ошибка: несовместимые типы

// Unary+ форсирует приведение к func* (positive hack):
auto g = +[]{};
g = []{};  // ok — g теперь указатель на функцию
\end{lstlisting}
\end{frame}

% ----------------------------------------------------------------
\begin{frame}[fragile]{Что и как захватывается}
\begin{lstlisting}
int g = 1;
void foo() {
  int x = 10; static int s = 2;

  auto lam = [=] { return x + s + g; };
  auto lam2 = [&] { return x + s + g; };
  auto lam3 = [&x, copy_s=s, copy_g=g] {
    return x + copy_s + copy_g;
  };
  x = 99;
  std::cout << lam() << std::endl; // 13
  std::cout << lam2() << std::endl; // 102
  std::cout << lam3() << std::endl; // 102
  s = 10;
  std::cout << lam() << std::endl;  // 21
  std::cout << lam2() << std::endl;  // 110
  std::cout << lam3() << std::endl;  // 102
}
// Правило: захватывается только локальный
// нестатический контекст. static и глобальные
// доступны напрямую и изменения в них всегда видны.
\end{lstlisting}
\end{frame}

% ----------------------------------------------------------------
\begin{frame}[fragile]{Провисшая ссылка}
\begin{lstlisting}
#include <iostream>

auto foo(int x = 10) {
  return [&]{return x;};
}

int main() {
  auto f = foo();
  std::cout << f(); // 32766
  return 0;
}
\end{lstlisting}
\end{frame}

% ----------------------------------------------------------------
\begin{frame}[fragile]{Взятие по ссылке в многопоточном коде}
\begin{lstlisting}
void byref() {
  std::vector<std::thread> threads;
  int i;
  for (i = 0; i < 10; i++) {
    threads.emplace_back([&i] {
       std::cout << i << ' '; });
  }
  for (auto &th : threads) { th.join(); }
  // 2 6 7 7 8 99  10 10 10
  std::cout << '\n';
}

void bycopy() {
  std::vector<std::thread> threads;
  for (int i = 0; i < 10; i++) {
    threads.emplace_back([i] {
        std::cout << i << ' '; });
  }
  for (auto &th : threads) { th.join(); }
  // 0 1 2 3 4 5 6 7 9 8
  std::cout << '\n';
}
\end{lstlisting}
\end{frame}



% ----------------------------------------------------------------
\begin{frame}[fragile]{mutable и захват с переименованием}
\begin{lstlisting}
// [a] → operator() const: нельзя менять копию
// [a]() mutable → без const: можно
auto counter = [a = 0]() mutable { return ++a; };
// counter() → 1, counter() → 2, ...

// Захват по const-ссылке: нужно переименование (C++14)
int x = 5;
auto lmd = [&rx = std::as_const(x), vx = x] {
    return vx + rx;  // rx: const int&
};
// lmd() → 10, но rx нельзя изменить

// Move-захват (C++14 init-capture)
auto ptr = std::make_unique<int>(42);
auto f = [p = std::move(ptr)]() { return *p; };
// ptr == nullptr после move
// f() == 42
\end{lstlisting}
\end{frame}

% ----------------------------------------------------------------
\begin{frame}[fragile]{Generic lambdas (C++14 / C++20)}
\begin{lstlisting}
// C++14: auto-параметр → шаблонный operator()
auto adder = [](auto x, auto y) { return x + y; };
adder(1, 2);      // int   → 3
adder(1.5, 2.5);  // double → 4.0

// C++20: явные шаблонные параметры + concepts
template <typename T, typename U>
concept Addable = requires(T t, U u) { t + u; };

auto typed = []<typename T, typename U>(T x, U y)
    requires Addable<T, U> { return x + y; };

typed(1, 2);       // ok
// typed("a", 1);  // ошибка: нет Addable<const char*,int>

// Доступ к типу T по имени внутри тела:
auto front = []<typename T>(std::vector<T>& v) -> T& {
    return v.front();
};
\end{lstlisting}
\end{frame}

% ----------------------------------------------------------------
\begin{frame}[fragile]{``Перегрузка'' лямбд (overload trick)}
\begin{lstlisting}
// Наследование от нескольких лямбд (C++17)
template <typename... F>
struct overload : F... { using F::operator()...; };
template<class... Ts>
overload(Ts...) -> overload<Ts...>;  // deduction guide

auto f = overload {
    [](int i)    { std::cout << "int "    << i; },
    [](double d) { std::cout << "double " << d; },
};
f(42); f(3.14);  // int 42 / double 3.14

// std::variant + overload ≈ паттерн-матчинг
std::variant<int, double, std::string> v = 3.14;
std::visit(overload {
    [](int i)         { std::cout << i; },
    [](double d)      { std::cout << d; },
    [](std::string& s){ std::cout << s; }
}, v);  // выведет: 3.14
\end{lstlisting}
\end{frame}

% ----------------------------------------------------------------
\begin{frame}[fragile]{std::function: стирание типов и рекурсия}
\begin{lstlisting}
// auto: нулевой overhead, конкретный тип лямбды
auto fast = [x = 5](int a) { return x + a; };

// std::function: type erasure, heap indirection
std::function<int(int)> flex = [x = 5](int a) {
    return x + a;
};  // внутри — void* + виртуальный вызов

// Рекурсивная лямбда требует стирания типа:
// auto fact = [&fact](int n){...}; // ошибка: fact ещё неизвестен
std::function<int(int)> fact = [&fact](int n) {
    return n <= 1 ? 1 : n * fact(n - 1);
};
fact(5);  // 120
// Внимание: fact захвачена по ссылке — не перемещать!

// Deducing this
auto fact = [](this auto&& self, int n) -> int {
    return n <= 1 ? 1 : n * self(n - 1);
};
fact(5);  // 120

\end{lstlisting}
\end{frame}

\end{document}
